%%%%%%%%%%%%%%%%%%%%%%%%%%%%%%%%%%%%%%%%%%%%%%%%%%%%%%%%%%%%%%%%%%%%%%%%
%% Two-page condensed CV --- Danilo Pianini                            %%
%% Self-contained: needs biber, but NO shell escape and NO network.    %%
%% The full CV lives in curriculum.tex.                                %%
%%%%%%%%%%%%%%%%%%%%%%%%%%%%%%%%%%%%%%%%%%%%%%%%%%%%%%%%%%%%%%%%%%%%%%%%
\documentclass[a4paper,9pt]{extarticle}

\usepackage[utf8]{inputenc}
\usepackage[T1]{fontenc}
\usepackage{textcomp}
\usepackage{xcolor}
\usepackage{microtype}
\usepackage{enumitem}
\usepackage{fancyhdr}
\usepackage[hyphens]{url}
\usepackage{pdftexcmds}
\usepackage{xstring}
\usepackage[
  backend=biber,
  style=ieee,
  sorting=ydnt,
  url=false,
  doi=false,
  isbn=false,
  maxbibnames=3,
  minbibnames=1,
  bibencoding=utf8
]{biblatex}
\usepackage[a4paper,margin=15mm,includefoot,footskip=12pt]{geometry}
\definecolor{darkblue}{rgb}{0.0,0.0,0.3}
\usepackage[hidelinks]{hyperref}
\hypersetup{colorlinks,breaklinks,linkcolor=darkblue,urlcolor=darkblue,
            anchorcolor=darkblue,citecolor=darkblue}

\addbibresource{bibliography.bib}

%%%%%%%%%%%%%%%%%%%%%%%%%%%%% Last modification date %%%%%%%%%%%%%%%%%%%%
% Reads the modification timestamp of this very source file, so that a
% rebuild without an edit does not change the printed date.
% \pdf@filemoddate yields D:YYYYMMDDhhmmss+ZZ'zz'
\makeatletter
\edef\cvmodraw{\pdf@filemoddate{\jobname.tex}}
\makeatother
\IfStrEq{\cvmodraw}{}
  {\newcommand{\cvmoddate}{\today}}
  {\newcommand{\cvmoddate}{%
     \StrMid{\cvmodraw}{3}{6}-\StrMid{\cvmodraw}{7}{8}-\StrMid{\cvmodraw}{9}{10}}}

%%%%%%%%%%%%%%%%%%%%%%%%%%%%% Bibliometrics %%%%%%%%%%%%%%%%%%%%%%%%%%%%%
% scholar-summary.tex is generated by scholar_scraper.rb during the
% curriculum.tex build; the \providecommand below is the offline fallback.
\IfFileExists{scholar-summary.tex}{\input{scholar-summary.tex}}{}
\providecommand{\scholarsummary}{3129 citations, h-index 28, i10-index 62
  (Google Scholar, 2026-09-15)}

%%%%%%%%%%%%%%%%%%%%%%%%%%%%%%% Layout %%%%%%%%%%%%%%%%%%%%%%%%%%%%%%%%%%
\setlength{\parindent}{0pt}
\setlength{\parskip}{0pt}
\linespread{0.96}\selectfont
\setlist{nosep,leftmargin=*,topsep=1pt,partopsep=0pt,parsep=0pt,itemsep=0pt}
\pagestyle{fancy}
\fancyhf{}\renewcommand{\headrulewidth}{0pt}
\fancyfoot[L]{\scriptsize\color{gray}Last updated: \cvmoddate}
\fancyfoot[C]{\scriptsize\color{gray}\thepage/2}
\fancyfoot[R]{\scriptsize\color{gray}Full CV:
  \href{http://www.danilopianini.org/}{danilopianini.org}}

% Full-width compact section heading
\newcommand{\cvsection}[1]{%
  \par\vspace{3.5pt}%
  {\normalsize\bfseries\scshape #1}\par\vspace{1pt}%
  \hrule height 0.6pt\vspace{2.5pt}\par}

% Entry with a right-flushed date; \hspace guarantees a gap before the date
\newcommand{\entry}[2]{\textbf{#1}\hspace{1em}\hfill{\small #2}\par}

% Bold the CV owner in the bibliography
\renewcommand*{\mkbibnamefamily}[1]{\ifstrequal{#1}{Pianini}{\textbf{#1}}{#1}}

\begin{document}

%%%%%%%%%%%%%%%%%%%%%%%%%%%%%%% Header %%%%%%%%%%%%%%%%%%%%%%%%%%%%%%%%%%
{\LARGE\bfseries Danilo Pianini}\hfill
{\small Italian citizen}\par
\vspace{1pt}\hrule height 1pt\vspace{2.5pt}

\textbf{Associate Professor} (\emph{Professore associato}) ---
Department of Computer Science and Engineering, Alma Mater Studiorum---Universit\`a di Bologna\par
\vspace{1pt}
{\small
Via dell'Universit\`a 50, 47522 Cesena (FC), Italy \quad$\cdot$\quad
+39 0547 33 88 20 \quad$\cdot$\quad
\href{mailto:danilo.pianini@unibo.it}{danilo.pianini@unibo.it} \quad$\cdot$\quad
\href{http://www.danilopianini.org/}{www.danilopianini.org}\par}

%%%%%%%%%%%%%%%%%%%%%%%%%% National habilitation %%%%%%%%%%%%%%%%%%%%%%%%
\cvsection{National scientific habilitation}
\setlength{\fboxsep}{4pt}%
\noindent\colorbox{black!7}{%
\begin{minipage}{\dimexpr\textwidth-2\fboxsep\relax}
\textbf{\large Abilitazione Scientifica Nazionale --- professore di I fascia (Full Professor)}\par
\vspace{1pt}
Sector \textbf{09/H1} --- Information processing systems.
Italian Ministry of University and Research.
Valid \textbf{2024-11-19} to \textbf{2036-11-19}
(\href{https://asn23.cineca.it/pubblico/miur/esito-abilitato/09\%252FH1/1/2}{decree}).
\end{minipage}}\par
\vspace{2pt}
{\small Also habilitated as \emph{professore di II fascia} in sector 01/B1 Informatics
(\href{https://asn18.cineca.it/pubblico/miur/esito-abilitato/01\%252FB1/2/5}{2020-11-23 to 2032-11-23})
and in sector 09/H1
(\href{https://asn16.cineca.it/pubblico/miur/esito-abilitato/09\%252FH1/2/5}{2018-07-26 to 2030-07-26}).\par}

%%%%%%%%%%%%%%%%%%%%%%%%%%% Research themes %%%%%%%%%%%%%%%%%%%%%%%%%%%%%
\cvsection{Research themes}
Engineering of pervasive computing: robust, coherent toolchains and methods leading to
adaptive, self-healing, and possibly evolving software ecosystems.
Design of novel programming paradigms and languages --- notably aggregate computing ---
that change how collective, large-scale, cyber-physical systems are modelled and built.

%%%%%%%%%%%%%%%%%%%%%%%%%%%%%% Education %%%%%%%%%%%%%%%%%%%%%%%%%%%%%%%%
\cvsection{Education}
\begin{itemize}
\item \entry{Ph.D.\ in Electronics, Computer Science and Telecommunications Engineering}{Universit\`a di Bologna, 2015}
      Thesis: \href{http://amsdottorato.unibo.it/7000/}{\emph{Engineering Complex Computational Ecosystems}}.
      Supervisor: Prof.\ Mirko Viroli.
\item \entry{M.Sc.\ in Computer Engineering --- 110/110 \emph{magna cum laude}}{Universit\`a di Bologna, 2011}
      Thesis: \emph{A Framework for the Simulation of Pervasive Services Ecosystems}.
      Supervisor: Prof.\ Mirko Viroli.
\item \entry{B.Sc.\ in Computer Engineering}{Universit\`a di Bologna, 2008}
      Thesis: \emph{From Swarm Intelligence to Self-Organising Coordination}.
      Supervisor: Prof.\ Andrea Omicini.
\end{itemize}

%%%%%%%%%%%%%%%%%%%%%%%%%%%% Academic career %%%%%%%%%%%%%%%%%%%%%%%%%%%%
\cvsection{Academic positions}
\begin{itemize}
\item \entry{Associate Professor, Universit\`a di Bologna}{2025-10 --- today}
\item \entry{Senior Assistant Professor (RTD-b), tenure track, Universit\`a di Bologna}{2022-10 --- 2025-10}
\item \entry{Junior Assistant Professor (RTD-a), Universit\`a di Bologna --- FluidWare PRIN project}{2019-12 --- 2022-10}
\item \entry{Research fellow (\emph{assegnista di ricerca}), Universit\`a di Bologna}{2015--2019, 2011}
      Aggregate programming platforms and toolchains; engineering collective adaptive systems.
\end{itemize}

%%%%%%%%%%%%%%%%%%%%%%%%%%%% Publications %%%%%%%%%%%%%%%%%%%%%%%%%%%%%%%
\cvsection{Publications and bibliometrics}
\textbf{122 peer-reviewed publications}: 37 journal articles, 63 papers in conference
proceedings, 22 books, chapters, and edited volumes.
\href{https://scholar.google.com/citations?user=0CLlt5oAAAAJ&hl=en}{Google Scholar}:
\textbf{\scholarsummary}.\par
\vspace{2pt}
\textit{Selected publications:}\par
\nocite{DBLP:journals/tosem/CasadeiAADPSTV25,JLAMP2019,TOMACS2018,
        BealIEEEComputer2015,PianiniSAC2015,PianiniJOS2013}
{\small\printbibliography[heading=none]}

%%%%%%%%%%%%%%%%%%%%%%%%%%%%%%% Awards %%%%%%%%%%%%%%%%%%%%%%%%%%%%%%%%%%
\cvsection{Awards}
\begin{itemize}
\item \entry{Best Student Paper Award --- IEEE ACSOS 2025, Tokyo, Japan}{2025}
      \emph{A Field-based Approach for Runtime Replanning in Swarm Robotics Missions}.
\item \entry{Best Paper Award --- 12th Int.\ Workshop on Engineering Multi-Agent Systems (EMAS), Auckland}{2024}
      \emph{On the external concurrency of current BDI frameworks for MAS}.
\item \entry{Best Paper Award --- IEEE SASO 2016, Augsburg, Germany}{2016}
      \emph{Self-Adaptation to Device Distribution Changes}; best among 20 accepted papers,
      invited for extension to ACM TAAS.
\item \entry{Best Poster Awards --- FSEN 2025 (V\"aster\r{a}s) and IEEE ACSOS 2024 (Aarhus)}{2024, 2025}
\end{itemize}

%%%%%%%%%%%%%%%%%%%%%%%% Funding and leadership %%%%%%%%%%%%%%%%%%%%%%%%%
\cvsection{Research leadership and funding}
\begin{itemize}
\item \entry{Scientific lead, \href{https://pslab-unibo.github.io/}{Pervasive Software Lab}, Universit\`a di Bologna}{2025 --- today}
      Research hub on pervasive computing hosting PhD students and research fellows;
      aggregate programming, digital twins, and their applications.
\item \entry{Unit leader and scientific PI --- WOOD4.0, Woodworking Machines for Industry 4.0}{2023-01 --- 2025-12}
      Emilia-Romagna regional project with SCM Group S.p.A.\ (leader) and Universit\`a di Bologna.
      \textbf{Budget of the UniBo unit: \texteuro{}300K}. CUP E69J22007520009.
\item \entry{Scientific lead of two industrial research contracts (SCM Group)}{2023 --- today}
      Remote software update platforms and cloud-edge continuum deployment for industrial machinery.
\item \entry{Research unit member --- PRIN national projects}{2019 --- today}
      CommonWears (PRIN 2022, 5 universities) and FluidWare (PRIN 2017, 4 universities);
      previously the EU SAPERE consortium. Sustained collaboration with Raytheon BBN
      Technologies (Dr.\ Jacob Beal), including a \$40K BBN--UniBo research contract.
\end{itemize}

%%%%%%%%%%%%%%%%%%%%%%%%%%%%% Supervision %%%%%%%%%%%%%%%%%%%%%%%%%%%%%%%
\cvsection{Supervision}
\textbf{5 PhD students:}
Angela Cortecchia (2024--), Davide Domini (2023--), Nicolas Farabegoli (2023--),
Martina Baiardi (2023--), Ruslan Shaiakhmetov (2022--2026).
Supervised student work has led to publications in top journals (e.g.\ \emph{Future
Generation Computer Systems}, 2024).\par
\vspace{1pt}
\textbf{67 graduating students:} 36 M.Sc.\ and 31 B.Sc.\ theses supervised, 2020--2026.

%%%%%%%%%%%%%%%%%%%%%%%%%%%%%%% Service %%%%%%%%%%%%%%%%%%%%%%%%%%%%%%%%%
\cvsection{Scientific community service}
\begin{itemize}
\item \entry{General Chair, \href{https://2026.acsos.org/}{IEEE ACSOS 2026}}{2026}
      The reference international conference on autonomic computing and self-organizing systems.
\item \entry{Program Committee Chair, IEEE ACSOS 2021}{2021}
      Further chairing roles: Track Chair MADTECC@PerCom 2024; Social Events Chair ACSOS 2023;
      General Chair ACSOS 2022 Italy; Special Event Chair ACSOS 2022; Track Chair NGPS@SAC 2019;
      Workshops and Tutorials Chair SASO 2018 and IEEE ICAC 2018; PC Chair ALP4IoT 2017.
\item \entry{Editorial roles}{2017 --- 2022}
      Academic editor, \emph{Scientific Programming}; associate blog editor, \emph{IEEE Software
      Blog}; topic board member, \emph{IJGI}; editor of the Springer volume \emph{The Future of
      Digital Democracy}; guest editor of the ACM TAAS ACSOS 2021 special issue.
\item \entry{Reviewing and committees}{2013 --- today}
      34 program committees of international conferences, reviewer for 22 international journals
      (incl.\ ACM TOSEM), 5 session chair roles, organizer of the Advanced School in Artificial
      Intelligence (Bertinoro, 2023), panelist at ACSOS 2022 and eCAS 2017.
\item \entry{IEEE Senior Member}{2026 --- today}
\end{itemize}

%%%%%%%%%%%%%%%%%%%%%%%%% Institutional roles %%%%%%%%%%%%%%%%%%%%%%%%%%%
\cvsection{Institutional roles (Universit\`a di Bologna)}
\begin{itemize}
\item \entry{Delegate of the Head of Department for Open Science, DISI}{2025-11 --- today}
\item \entry{Member of the academic board, PhD programme in Computer Science and Engineering}{2025-05 --- today}
\item \entry{President of the Internationalisation Commission (B.Sc.\ Computer Science and Engineering)}{2025-03 --- today}
\item \entry{FAIR Champion --- promoting FAIR principles across the University research community}{2024-12 --- today}
\item \entry{Dual Career Program academic tutor (B.Sc.\ and M.Sc.\ Computer Science and Engineering)}{2022-03 --- today}
\end{itemize}

%%%%%%%%%%%%%%%%%%%%%%%%%%%%%% Teaching %%%%%%%%%%%%%%%%%%%%%%%%%%%%%%%%%
\cvsection{Teaching}
\begin{itemize}
\item \entry{Software Process Engineering --- M.Sc.\ in Computer Science and Engineering}{2023/24 --- 2025/26}
      Course responsible. 3 CFU, 30 hours, $\sim$25 students/year.
\item \entry{Laboratory of Software Systems --- M.Sc.\ in Computer Science and Engineering}{2020/21 --- 2022/23}
      Course responsible. 3 CFU, 30 hours, $\sim$30 students/year.
\item \entry{Software Design and Development --- B.Sc.\ in Computer Systems Technologies}{2025/26}
      Course responsible. 3 CFU, 30 hours.
\item \entry{Further teaching}{2011 --- today}
      20 learning modules, 2 PhD courses (DevOps for scientific research; developing and sharing
      research software), and 8 tutoring assignments.
      Student satisfaction on responsible courses consistently between 75\% and 100\%.
\end{itemize}

%%%%%%%%%%%%%%%%%%%%%%%%%%%%%% Software %%%%%%%%%%%%%%%%%%%%%%%%%%%%%%%%%
\cvsection{Scientific software leadership}
\begin{itemize}
\item \entry{\href{http://alchemistsimulator.github.io/}{Alchemist} --- lead designer and main developer}{2010 --- today}
      Simulator joining agent-based modelling with stochastic simulation algorithms; widely
      adopted for pervasive computing, collective systems, and computational biology scenarios.
\item \entry{\href{http://protelis.org/}{Protelis} --- lead designer and main developer}{2014 --- today}
      Aggregate programming language making heterogeneous networked systems as easy to program
      as single machines.
\item \entry{\href{http://collektive.github.io/}{Collektive} --- lead designer and main developer}{2023 --- today}
      Kotlin multiplatform implementation of aggregate computing, embedding its semantics into
      Kotlin through a compiler plugin.
\end{itemize}

%%%%%%%%%%%%%%%%%%%%%%%% International experience %%%%%%%%%%%%%%%%%%%%%%%
\cvsection{International research visits}
\begin{itemize}
\item \entry{University of Iowa, Iowa City, IA, USA --- Visiting Researcher (Dr.\ Jacob Beal)}{2016, 2014}
\item \entry{Raytheon BBN Technologies, Cambridge, MA, USA --- Visiting Researcher (Dr.\ Jacob Beal)}{2014}
      Design and realisation of the Protelis programming language.
\item \entry{Johannes Kepler Universit\"at, Linz, Austria --- Visiting Researcher (Prof.\ Alois Ferscha)}{2013}
\item \entry{Florida Institute of Technology, Melbourne, FL, USA --- Visiting Researcher (Dr.\ Ronaldo Menezes)}{2009}
\end{itemize}

%%%%%%%%%%%%%%%%%%%% Dissemination and knowledge transfer %%%%%%%%%%%%%%%
\cvsection{Dissemination and knowledge transfer}
\begin{itemize}
\item \entry{Talks}{2009 --- today}
      21 talks at international conferences and workshops, including an invited Doctoral
      Symposium talk at \href{https://2024.acsos.org/}{IEEE ACSOS 2024}
      (\emph{Reproducible experiments, or didn't happen}), plus institutional seminars and
      public-engagement lectures.
\item \entry{Open-source contributions}{2008 --- today}
      Contributor to \href{https://jakta-bdi.github.io/}{JaKtA} (Kotlin BDI agents) and to
      \href{https://github.com/renovatebot/renovate}{Mend Renovate} (Gradle version catalogs
      support); author of a dozen widely-used Gradle and JVM developer tools published on
      Maven Central.
\item \entry{Industrial consulting}{2016 --- today}
      Peer Network (DevOps, build automation, continuous integration, Kotlin; 2021, 2024)
      and twinlogix (code generators; 2016, 2021).
\item \entry{Extra-institutional teaching}{2016 --- 2019}
      \href{https://www.bbs.unibo.eu/hp/}{Bologna Business School}, master course on IoT
      software production; \href{http://www.formart.it/}{FORMart}, four higher technical
      education courses (30 hours each) on IoT, programming, and computing systems.
\item \entry{Board member, \href{https://www.distrettodellinformaticaromagnolo.org/}{Distretto dell'Informatica Romagnolo} (DIR)}{2025-04 --- today}
      Non-profit community building networks between software companies and educational
      institutions in the Romagna area.
\end{itemize}

\end{document}
